% Nouveau EGA V (cote 25): archive cote 25, site page 164 (PDF page 165).
% English translation of the checked French annotated master.
% AI-translated by GPT-6 Astra (OpenAI, ChatGPT Pro mode).
\scanpage{25/164}
\mspage{\hand{7}}

k-isomorphic to $k[T]/\var{\over{T}{(}}{(}{writing over}{ega25_25_164-v0001}T^n)$ . If $d=1$,\var{\inserted[typed]{x rat. over k,}}{x rat. over k,}{insertion}{ega25_25_164-v0002} the notion of an
\var{\xout{ordinary}}{}{typewriter strike-out}{ega25_25_164-v0003} elementary singular point of multiplicity n corresponds to the classical
terminology \var{\xout{``curve}}{}{typewriter strike-out}{ega25_25_164-v0004}``ordinary singular point with n distinct tangent\unsure{s|z}''.
\var{\margin{(}}{\placedmargin{(}}{marginal note placed}{ega25_25_164-v0005}(We shall \var{\inserted{already}}{already}{insertion}{ega25_25_164-v0006} have made explicit in 15.2. that if x is a point of X rat over k, then
the notion of an elementary singularity of mult. n, in the absolute or
relative sense, \var{\deleted{it}}{}{deletion}{ega25_25_164-v0007} is the same; this also remains valid if k(x)
is a finite separable extension of k---probably, moreover, ``finite''
is unnecessary---and this fact deserves to be inserted as a corollary \var{\xout{if one succeeds}}{}{typewriter strike-out}{ega25_25_164-v0008}.
or as a proposition\hand{)}.

NB I had included in par.20 a no. on the blowing up of a
prescheme X along a \var{\inserted[typed]{closed}}{closed}{insertion}{ega25_25_164-v0009} subprescheme Y such that the augmentation
$\underline{O}_X \hand{\longrightarrow} \underline{O}_Y$ is elementary of mult.n (to which I alluded in
15.6.). I point out that this is a short and welcome no., whose
only ingredients are the general smoothness results of par.17 and
the definition \var{\inserted[typed]{at the beginning of}}{at the beginning of}{insertion}{ega25_25_164-v0010} 17.15.\unsure{2}. \var{\inserted[typed]{(and even this last one could be avoided).}}{(and even this last one could be avoided).}{insertion}{ega25_25_164-v0011} Nothing would prevent incorporating these results
here, \var{\xout{by in}}{}{typewriter strike-out}{ega25_25_164-v0012} for example immediately after Definition
\var{\xout{17}}{}{typewriter strike-out}{ega25_25_164-v0013}15.1. above; then \var{\inserted[typed]{the ``geometric'' range of 15.2. and the sequence}}{the ``geometric'' range of 15.2. and the sequence}{insertion}{ega25_25_164-v0014} 15.\unsure{3|2}. to 15.7. could be separated
from these \var{\xout{calculations}}{}{typewriter strike-out}{ega25_25_164-v0016} results by grouping them in a No. 17.16, \var{\xout{under a title}}{}{typewriter strike-out}{ega25_25_164-v0017}
\var{\xout{such as}}{}{typewriter strike-out}{ega25_25_164-v0018} In any case, I realize as I write that the
\var{\margin{(}}{\placedmargin{(}}{marginal note placed}{ega25_25_164-v0019}new no. 15 planned on smooth forms breaks up into \var{\inserted[typed]{at least}}{at least}{insertion}{ega25_25_164-v0020} two fairly
distinct and independent numbers, one containing general sorites for
the ``elementary of mult n'' things, with arbitrary n, the other \var{\deleted{which}}{}{deletion}{ega25_25_164-v0021} contain\var{\over{s}{ing}}{ing}{writing over}{ega25_25_164-v0022} the
characterization of smooth \uline{quadratic} forms, and which takes from
the preceding one only no. 15.1., i.e. practically the definition of the smoothness
of a form.

\mspage{\hand{164}}
\tnote{25-164-1}{Rat., mult., par. and no. mean rational, multiplicity, paragraph and numbered subsection. The source's disctinctes and incopporer are translated without creating English misspellings. The canceled doubtful \unsure{elle} is translated it. The word probably follows the terminology standard for sans doute; the author's qualification concerning finite is retained.}
\tnote{25-164-2}{Range renders gamme in the planning passage, and sorites retains the supplied glossary term. The canceled isolated t after à has no English counterpart. The last writing-over changes French contient to contenant; English represents the same change by contains to containing, with its suffix variant retained. All proposed section numbers and their doubts remain as written.}
